\chapter{话题漂移感知的对话历史剪枝}
\label{chap:hpprune}

本章提出并详述\textbf{话题漂移感知的对话历史剪枝器}（History Pruner, HP-Pruner），作为本文三阶段查询重写框架的第一环。第~\ref{sec:hp_motivation}~节从实证角度揭示"长历史即高噪声"现象；第~\ref{sec:hp_model}~节给出 HP-Pruner 的模型设计；第~\ref{sec:hp_data}~节介绍弱监督训练数据的自动构造方案；第~\ref{sec:hp_train}~节给出训练目标与推理策略；第~\ref{sec:hp_exp}~节报告在 TopiOCQA 与 QReCC 上的消融与分析实验。

\section{动机与问题分析}
\label{sec:hp_motivation}

\subsection{长历史带来的三类负面影响}

在多轮对话检索场景中，把全部历史 $H_i$ 简单拼接作为重写器输入会带来三方面负面影响：

\textbf{噪声注入}：当 $H_i$ 中存在与 $q_i$ 无关或仅弱相关的历史轮次时，这些轮次中的实体、动词、限定词容易被重写器误引入最终重写，造成"幻觉式扩展"。

\textbf{计算成本}：大模型重写器的推理开销与输入长度近似成线性关系，把全部历史塞入上下文窗口会放大推理延迟 2--5 倍。

\textbf{话题切换误导}：当对话发生硬话题切换（后一轮与前数轮完全无关）时，重写器若仍试图"统一"两个话题，极易产出语义扭曲的混合查询，检索器无法召回。

\subsection{实证观测}

我们在 TopiOCQA 的有切换标注子集上观测到：若直接把完整历史拼接后用 BGE-large 作为检索器，Recall@5 仅为 15.2；切换到稍强的 GTE-ModernBERT 也只能得到 29.4；当用本文 §\ref{chap:prw} 提出的渐进式重写器对输入进行整理后，同样的检索器组合达到 54.1 和 58.9，相差 30+ 个百分点（详见表~\ref{tab:main_r5}）。这组对照直接证明：\textbf{历史轮次并非越多越好，有针对性的剪枝与重写能带来显著增益}。

\section{模型设计}
\label{sec:hp_model}

\subsection{三分类轮次级相关性判别器}

本文将历史轮次与当前查询的关系划分为三类：
\begin{itemize}
    \item \textbf{强相关（label 2）}：同一话题且包含重要实体或限定词；
    \item \textbf{弱相关（label 1）}：同一话题但不提供额外信息；
    \item \textbf{无关（label 0）}：已切换到新话题，与当前查询无语义关联。
\end{itemize}

形式化地，对每个 $(q_j, a_j) \in H_i$ 轮次，分类器预测：
\begin{equation}
s_{j,i} = f_\phi\big(\text{[CLS]} \oplus q_i \oplus \text{[SEP]} \oplus (q_j \oplus a_j)\big) \in \{0,1,2\},
\end{equation}
其中 $f_\phi$ 以 ModernBERT-base~\citep{warner2024modernbert} 为骨干编码器，其后接两层 MLP 与 softmax 分类头。选择 ModernBERT 的动机是其更长的上下文窗口（8192 tokens）与更低的推理成本，使得对长历史的轮次级扫描变得实用。

\subsection{剪枝策略}

在推理阶段，按轮次顺序从最早到最近依次打分，应用阈值 $\tau$ 决定保留与否：
\begin{equation}
\tilde H_i = \{(q_j, a_j)\mid s_{j,i} \geq \tau\}.
\end{equation}
默认 $\tau=1$（保留强相关与弱相关）。为平衡上下文完整性与压缩比，额外保留紧邻当前轮次的前一轮（"last-turn guard"），以防分类器误杀重要上下文。

\section{弱监督训练数据构造}
\label{sec:hp_data}

获取大规模轮次级相关性标签成本高昂。本文提出一套\textbf{自动弱监督数据生成方案}，基于两类信号：

\subsection{基于话题切换标注的信号}

TopiOCQA 提供了每轮是否发生话题切换的标注。利用该信号可自动导出：若轮次 $j$ 与当前轮次 $i$ 之间发生过话题切换，则标记 $(q_j, a_j)$ 对 $q_i$ 为\textbf{无关}；否则默认为\textbf{弱相关}。

\subsection{基于 QReCC 人工重写的信号}

QReCC 为每轮提供人工去上下文化的重写 $q_i^{\text{gold}}$。利用如下启发式可区分强/弱相关：
\begin{equation}
\text{label}(j,i)=\begin{cases}
\text{强相关 (2)} & \text{if } \exists\, e \in \text{Entities}(q_j, a_j),\ e \in q_i^{\text{gold}}\\
\text{弱相关 (1)} & \text{otherwise, if } j \text{ 与 } i \text{ 同一话题}\\
\text{无关 (0)} & \text{if } j \text{ 与 } i \text{ 跨话题}
\end{cases}
\end{equation}
其中 Entities 用命名实体识别器 spaCy 自动抽取。

\subsection{合成困难负例}

为强化分类器对"跨话题但表面词汇相似"困难样例的判别能力，本文构造\textbf{合成困难负例}：从不同对话中随机抽取话题不同但共享命名实体的两条轮次配对，标记为 \textbf{无关}。最终将 TopiOCQA 话题切换信号与 QReCC 人工重写词级对比启发式合并，得到 94\,422 条轮次-查询对，三类标签（无关 / 弱相关 / 强相关）分布为 36\%\,/\,35\%\,/\,29\%，按 95:5 划分训练/验证。

\section{训练目标与实现细节}
\label{sec:hp_train}

采用加权交叉熵损失：
\begin{equation}
\mathcal{L}_{\text{HP}} = -\frac{1}{|\mathcal{B}|}\sum_{(j,i)\in \mathcal{B}} \sum_{c=0}^{2} w_c \cdot \mathbb{1}[y_{j,i}=c] \log p_\phi(c \mid q_i, q_j, a_j),
\end{equation}
其中类别权重 $w_c$ 按训练集中频率的倒数设置，以缓解类别不均衡带来的偏置。优化器 AdamW，学习率 $2\!\times\!10^{-5}$，batch size 32，warmup 比例 3\%，权重衰减 0.01，最大序列长度 512，训练 3 轮，在单张 A100 上约 20 分钟完成训练（共 8\,412 步）。

\section{实验与分析}
\label{sec:hp_exp}

\subsection{分类器本身性能}

表~\ref{tab:hp_classifier} 报告三分类任务在 4\,721 条留出验证集上的准确率、宏 F1 与加权 F1。整体准确率达到 80.1\%，宏 F1 为 0.797，表明 ModernBERT-base 规模的轻量编码器即可稳定刻画轮次级相关性。

\begin{table}[h]
\centering
\bicaption{HP-Pruner 轮次级三分类性能（留出验证集，$N=4\,721$）}{Turn-level three-way classification performance of HP-Pruner}
\label{tab:hp_classifier}
\begin{tabular}{lccc}
\toprule
指标 & Accuracy & Macro-F1 & Weighted-F1 \\
\midrule
HP-Pruner (ModernBERT-base) & 80.13 & 0.797 & 0.799 \\
\bottomrule
\end{tabular}
\end{table}

\subsection{三类标签的混淆模式}

进一步分析分类器在三类标签上的表现。"无关"类别（即硬话题切换后）由于来自 TopiOCQA 的话题段落标注，边界清晰，识别准确率最高；"弱相关"与"强相关"之间存在一定混淆，这与人工标注的本质不确定性一致——当前轮次可能同时依赖多轮上下文，硬划为二选一本身就带有主观性。从实际剪枝效果看，将"弱相关"与"强相关"合并保留（即 $\tau=1$）的策略可以很好地回避这一歧义。

\subsection{阈值敏感性}

在推理时扫描阈值 $\tau\in\{0,1,2\}$：$\tau=0$ 相当于保留全部历史（退化为 No-Rewrite 输入），$\tau=2$ 会过度激进地丢弃上下文导致下游重写器信息不足；$\tau=1$（保留"弱相关"与"强相关"）在验证集上取得准确率与压缩比的最佳权衡。错误分析显示，"长回复中包含关键实体但话题已切换"的样例最难判别，这类样例约占错误案例的三分之一，未来可通过更细粒度的 span 级相关性建模进一步改进。

\section{本章小结}

本章提出了话题漂移感知的轮次级三分类剪枝器 HP-Pruner，设计了基于话题切换标注与实体重叠信号的弱监督训练方案，并在 TopiOCQA 与 QReCC 上实证了其对下游重写与检索性能的显著增益。该模块将作为本文三阶段框架的第一环，为后续的渐进式重写提供干净、紧凑的上下文。
